\chapter{Anexos}

\noindent Los anexos reúnen detalles operativos que complementan la explicación principal: estructura del repositorio, configuración del entorno, comandos reproducibles y parámetros de entrenamiento.

\section{Estructura del repositorio}
\label{anx:repositorio}

El código fuente completo se encuentra en el repositorio público del proyecto:\\[0.4em]
\begin{center}
\url{https://github.com/STOSKR/TFM-CSFlipper}
\end{center}

La Figura~\ref{fig:estructura-repositorio} resume la organización. Se muestran solo los directorios que ayudan a localizar cada parte del sistema, mientras que los ficheros y resultados generados se consultan directamente en el repositorio.

\begin{figure}[htbp]
\centering
\fbox{%
\begin{minipage}{0.78\textwidth}
\small
\begin{tabular}{@{}ll}
\texttt{TFM-CSFlipper/} & \\
\quad \texttt{apps/} & comandos, extractores e interfaz local \\
\quad \texttt{packages/} & lógica reutilizable del sistema \\
\quad \texttt{tests/} & pruebas unitarias y de integración \\
\quad \texttt{data/} & datasets y resultados derivados \\
\quad \texttt{docs/} & documentación técnica \\
\quad \texttt{TFM/} & memoria, figuras y bibliografía
\end{tabular}
\end{minipage}}
\caption{Estructura del repositorio.}
\label{fig:estructura-repositorio}
\end{figure}

Los directorios principales cumplen estas funciones:

\begin{itemize}
  \item \texttt{apps}: contiene los extractores, los comandos de consola y la interfaz local.
  \item \texttt{packages}: agrupa los contratos, la persistencia, la predicción, el simulador, el entorno \acroref{MARL}{MARL} y el procesamiento visual.
  \item \texttt{tests}: reúne las pruebas unitarias y de integración que verifican reglas económicas, conectores y flujos principales.
  \item \texttt{data}: conserva datasets derivados, informes y resultados que no forman parte del código fuente.
  \item \texttt{docs}: recoge documentación técnica y decisiones de desarrollo.
  \item \texttt{TFM}: contiene el código \LaTeX{}, las figuras y la bibliografía de esta memoria.
\end{itemize}

\section{Configuración del entorno}

El entorno se define en \texttt{pyproject.toml}. El proyecto usa Python, Supabase/PostgreSQL, Pytest, Ruff, Mypy y dependencias opcionales para \acroref{MARL}{MARL} como Ray/\acroref{RLlib}{RLlib}, PettingZoo, Gymnasium y PyTorch. Para desarrollo local existe un \texttt{docker-compose.yml} con PostgreSQL.

Comandos representativos:

\begin{verbatim}
python -m pytest tests/unit
python -m ruff check
python -m mypy apps packages tests/unit
python -m apps.cli.auto_scrape_loop --once
python -m apps.cli.render_stream_scrape
python -m apps.cli.web_dashboard_server
\end{verbatim}

\section{Comandos de datasets y modelos}

Los comandos de datos y modelos separan adquisición, construcción de dataset, entrenamiento e inferencia. Esta separación permite repetir una fase sin rehacer todo el flujo.

\begin{verbatim}
python -m apps.cli.build_supervised_dataset
python -m apps.cli.analyze_supervised_coverage
python -m apps.cli.train_supervised_model
python -m apps.cli.build_trading_dataset
python -m apps.cli.retrain_trading_model
python -m apps.cli.score_live_opportunities
python -m apps.cli.train_marl_ctde
python -m apps.cli.run_marl_experiment_matrix --help
\end{verbatim}

Estos comandos son puntos de entrada, no una receta para ejecutar todos los experimentos sin revisión. Los procesos que consultan o persisten datos requieren la configuración local y las credenciales correspondientes; la matriz multiagente, además, recibe un plan y un directorio de salida. Los argumentos concretos de cada ejecución y sus resultados se conservan como artefactos, según se describe al final de este anexo.

\section{Parámetros de entrenamiento}

Los experimentos supervisados usan particiones temporales, calibración de probabilidades y selección por precisión operativa en umbrales altos. Las exploraciones comparan un modelo de referencia que predice la clase mayoritaria, regresión logística, \emph{random forest} e \emph{histogram gradient boosting}. El entrenamiento \acroref{MARL}{MARL} usa \acroref{RLlib}{RLlib} con un entorno multiagente y tres políticas: \textit{Scout}, \textit{Trader} y \textit{Portfolio}.

\begin{table}[h]
\centering
\begin{tabularx}{\textwidth}{>{\raggedright\arraybackslash}p{0.18\textwidth} >{\raggedright\arraybackslash}X}
\toprule
Bloque & Parámetros relevantes \\
\midrule
Supervisado & Partición temporal, umbral operativo, calibración, modelo candidato y columnas excluidas para evitar información futura. \\
Compraventa & Dirección de la operación, horizonte, comisiones, tipo de cambio, ventana histórica y umbral de beneficio neto. \\
Simulación & Saldo inicial, tamaño máximo de posición, \textit{trade hold}, comisiones y límites de exposición. \\
\acroref{MARL}{MARL} & Políticas \textit{Scout}/\textit{Trader}/\textit{Portfolio}, algoritmo \acroref{PPO}{PPO}, $\gamma=0{,}99$, recompensa, restricciones y métricas de cartera. \\
Matriz \acroref{MARL}{MARL} & 1.000~EUR iniciales, 100 iteraciones máximas, ocho episodios por iteración, parada tras 20 iteraciones sin mejora y semillas 7, 19 y 31. \\
\bottomrule
\end{tabularx}
\caption{Configuración experimental.}
\end{table}

Los valores de la última fila corresponden a la matriz comunicada en el Capítulo~7. No constituyen parámetros universales del modelo: se registran para que la comparación entre configuraciones pueda repetirse en las mismas condiciones.

\section{Modelo de datos}

El modelo operativo actual gira alrededor de dos tablas principales: \texttt{market\_items} y \texttt{market\_history\_points}. La primera representa artículos identificados por sus atributos y su estado actual por plataforma. La segunda almacena series temporales en formato largo, con una fila por artículo, plataforma, fecha y métrica. La separación entre \texttt{market\_items} y \texttt{market\_history\_points} evita mezclar métricas distintas y permite añadir nuevas señales sin alterar la forma general del histórico.

El repositorio conserva además un esquema canónico de migración, distinto del flujo experimental operativo actual. Incluye \texttt{assets}, \texttt{platforms}, \texttt{market\_observations}, \texttt{outbox\_events}, \texttt{predictions}, \texttt{risk\_profiles}, \texttt{votes}, \texttt{investment\_decisions} y \texttt{simulated\_positions}. Su objetivo es mejorar la trazabilidad, la auditoría y la separación entre observación, predicción y decisión. No se utiliza como fuente primaria de los conjuntos descritos en el Capítulo~7.

\section{Artefactos de reproducibilidad}

Los conjuntos derivados se guardan en \texttt{data/datasets}, incluidos los conjuntos \texttt{trading\_profit\_v1} y \texttt{trading\_profit\_v2}. Los modelos y las salidas de entrenamiento se almacenan en \texttt{model-runs}; los cortes temporales, informes y figuras de exploración se conservan en \texttt{data/experiments} y \texttt{docs/experiments}. Esta separación evita mezclar los artefactos generados con el código fuente y permite recuperar el contexto de una ejecución concreta.

Para que un experimento aporte evidencia reproducible, debe registrar como mínimo: configuración usada, periodo temporal, número de artículos, tamaño de cada partición, proporción de operaciones rentables, modelo o política evaluada, métricas principales y limitaciones observadas.
